\section{Implementation and experimental protocols}
\label{app:relu-details}

\subsection{aSRR implementation settings}
\label{app:asrr-construction}

We use a denominator floor of $10^{-12}$ for both column normalization
and Frobenius normalization. At each site, we divide the neurons into
equal contiguous blocks. The block count is the largest divisor of the
site width up to $16$.

We optimize the four rotation fractions for 144 simultaneous
perturbation updates and average the last 24 iterates. Each
classification update uses 64 development episodes and compares
refits at Adam rates 0.003 and 0.03 after 12, 75 and 300 steps.
The two perturbations and the signed basis use the same episodes
and starting values. Final selection evaluates the full development
sample with the numerical refit settings below.

\paragraph*{Alignment.}
The $\operatorname{Align}$ step in Algorithm~\ref{alg:asrr} uses
neuron permutations within each site~\cite[Sec.~2.1]{zhou2023permutation}.
For weight matrices $X,Y$, let $P(X,Y)$ be the within-site permutation
maximizing $\langle XP,Y\rangle_F$, computed by the Hungarian algorithm.
Write $s_\tau=\arg\min_s\ell_{\tau s}$ and
$\tau_*=\arg\min_\tau\ell_{\tau s_\tau}$ for the reference indices. The
routine returns the aligned means
\[
    \begin{aligned}
        \widetilde{\mathbf W}_\tau
            &=\operatorname{mean}_s\!
                \left[\mathbf X_{\tau s}P(\mathbf X_{\tau s},
                    \mathbf X_{\tau s_\tau})\right],\\
        \overline{\mathbf W}_\tau
            &=\widetilde{\mathbf W}_\tau
                P(\widetilde{\mathbf W}_\tau,\mathbf X_{\tau_*s_{\tau_*}}).
    \end{aligned}
\]

\subsection{Transfer benchmark protocol}
\label{app:transfer-details}

\subsubsection{Cross-correlation regression}

For each layout, we draw 96 source, 48 development and 160 test
functions. Each source function has 96 Gaussian input vectors and
independent response noise with standard deviation 0.05. We use 80
examples for each least-squares source fit. For $\mathbf K$ and
$\mathbf Q_1$, we split the 96 examples into training and validation
halves.

Each development or test function has 96 noiseless evaluation
examples. For aSRR we select one rotation using both training sample
sizes and choose a ridge strength for each sample size. Paired
intervals use the six layout means. Yeh and MSR-FC differentiate
through source ridge refits.
Meta-SP uses 80 training examples per source function and alternates
updates of size $0.9/L$ on the source loss with SVD projections to rank
54 over the given connections. Here $L$ is the largest eigenvalue
across the normalized Gram blocks. We select its checkpoint and ridge
strength on the same development functions.

\subsubsection{Classification architecture and data splits}

The classification network uses one hidden layer with coordinatewise
ReLU $\rho(t)=\max\{t,0\}$,
\[
    f(\mathbf x)=\mathbf W_2^{\mathsf T}
        \rho(\mathbf W_1^{\mathsf T}\mathbf x+\mathbf b_1)+\mathbf b_2,
\]
and sharing applies to the permitted connections in $\mathbf W_1$.
A site consists of neurons that receive the same input coordinates.

Toroidal MNIST and Aug-Omniglot use 64 hidden neurons. We divide the 45
MNIST digit pairs into 25 source, 10 development and 10 test pairs,
splitting original images before applying a random circular shift and
a fixed pixel permutation. We draw 2,000 two-way test episodes.
Aug-Omniglot uses 48 source tasks, 192 development episodes and 768
test episodes with disjoint character pools. Transformed images in
validation and test use random resized crops, reflections and rotations up
to $30^\circ$.

FC100 and Aug-miniImageNet use 16 patches with eight neurons per patch
and mean-pooled $4\times4$ RGB inputs. FC100 uses 16 source tasks,
384 development episodes and 2,048 test episodes. Aug-miniImageNet
uses the 64/16/20 class split of \citet{ravi2017optimization} and the
same transformations as Aug-Omniglot. Each of 64 source tasks supplies
64 clean training images and 30 validation images per class. Each
validation image has a clean and a transformed view. We draw 2,048
five-way test episodes with 15 test images per class.

WISDM uses 18 activities and nonoverlapping 64-sample windows of the
three smartwatch acceleration axes and their magnitude. The network
has 32 hidden neurons across four sites. We use 27 source, 7
development and 12 test participants with sufficient windows per
activity. The 768 test episodes use early windows for training and
later windows for testing.

UCI HAR uses the original features and widths 561--16--6. We retain
every second window within each subject and activity. Each of four
source tasks pools four participants and uses 366 training windows
and two separate validation samples of 78 windows for selecting Dense
fits and building the shared structure, respectively. Five participants supply 96
development episodes and nine supply 360 test episodes. Each test
episode uses three training and five test windows per activity.

\subsubsection{Source signals for classification}

We retain three Dense fits per source task on FC100 and four on the
other classification datasets. On Aug-Omniglot and Aug-miniImageNet,
the two validation samples contain the same images before and after
augmentation. For Equation~\eqref{eq:asrr-moments}, we set
$\varepsilon=0.03\max\{\operatorname{tr}(\sum_v\pi_v\Sigma_v)/d,10^{-8}\}$.

For each of eight training draws per source task at the target sample
size, we fit a reference network for 15 Adam steps at rate 0.01. Its
basis uses signed groups from $\mathbf A$ at the same rank, with at
most four equal blocks per site. We compute the training gradient and
the validation gradient for each version $v$ at that fit.

For a given block, $\pi_v$ weights validation sample $v$ in both the
gradient and class statistics. Let $V$ be the number of validation
samples and let $v_*$ index the one with the highest mean loss at the
reference fits. Set
$\omega_v=1/(2V)+\mathbf1_{\{v=v_*\}}/2$. At each site with at least
$V$ blocks, assign one block to each version. Assign the $b$th
block thereafter to the version maximizing $b\omega_v-n_v$, where $n_v$ is
the number of blocks already assigned to that version. We use
\[
    \pi_v=
    \begin{cases}
        \mathbf1_{\{v\text{ assigned to the block}\}}, & \text{at least $V$ blocks at the site},\\
        \omega_v, & \text{fewer than $V$ blocks at the site}.
    \end{cases}
\]

\subsubsection{Classification refits and evaluation}

For the final aSRR refit on MNIST, Aug-Omniglot, FC100 and WISDM,
we compare Adam rates 0.001, 0.003, 0.01 and 0.03 with 25, 75, 150
or 300 steps. We average accuracy over two starts. We break ties in
the smallest improvement across validation samples by mean accuracy
and then by validation cross-entropy. HAR uses the same rates,
activation quantiles 0.1, 0.5
and 0.9, and up to 1,000 steps. We retain three settings by validation
cross-entropy and compare ensembles of 1, 4 and 16 predictions.

For Aug-miniImageNet, we refine source checkpoints and refits on 1,152
development episodes. For the final aSRR refit we compare rates
0.02, 0.03, 0.045 and 0.06 with 5, 8, 10, 12, 16, 25 or 75 steps.
Development selects 8 steps at rate 0.045. The selected refits use
10 steps at rate 0.045 for MSR-FC, 1,200 at 0.005 for Yeh (soft) and
300 at 0.01 for Dense. For WSR, LUT-Q and Yeh (hard),
we first select three indexators per method by mean validation
accuracy and then refine the Adam rate and duration.

All methods refit on the same test episodes. Accuracy averages two
fresh starts, with HAR using the selected ensemble. Paired intervals
use participant means on HAR and WISDM, digit-pair means on toroidal
MNIST and episodes conditional on the test class pools on FC100,
Aug-Omniglot and Aug-miniImageNet.

\subsubsection{Classification baselines and controls}

On the image datasets and WISDM, WSR and LUT-Q propose indexators from
the best Dense fit on each of the eight source tasks with lowest
validation loss. WSR uses the isotonic proximal map with
rewinding~\citep{shakerinava2024weight}. We set the sparsity penalty to
zero and select the proximal strength at each update to match the code
budget. LUT-Q uses quantized forward passes and straight-through
gradients. Each update reassigns weights to their nearest entries and
recomputes dictionary means~\citep{cardinaux2018iteratively}. Both
methods transfer indexators that occupy every code. The HAR LUT-Q
dictionary uses signed copies of each coefficient and a fixed zero.

Yeh's method~\citep{yeh2022equivariance} optimizes validation loss
through source refits. On image datasets and WISDM, we differentiate
through three inner updates and use entropy-regularized rectangular
assignments per block. Hard decoding maximizes total assignment
probability subject to occupying every code. The soft variant
transfers an orthonormal basis of the learned span. On Aug-miniImageNet,
the hard variant compares warm and random initializations, rates
0.003 and 0.012, and checkpoints from 800 through 6,400 source
updates. HAR uses signed assignments with a fixed zero, a 20-term
Neumann hypergradient approximation and rowwise maximum-probability
hard decoding. The selected HAR soft variant uses entropy and
nuclear-norm penalties.

MSR-FC~\citep{zhou2021meta} learns a first-layer reparameterization by
differentiating validation loss through source refits. On image
datasets and WISDM, it learns separable input and output factors. At
transfer, we orthonormalize the input factors and absorb the invertible
output factors into the coefficients. HAR uses an unconstrained basis.
On Aug-Omniglot, development selects four input factors over one. Each
$784\times32$ factor serves 16 neurons with separate coefficients.
The original augmentation experiments use reparameterized convolutional
networks with meta-learned initial values and SGD rates. Our
Aug-miniImageNet source objective uses both validation views.

\phantomsection\label{app:spectral-controls}
For each neuron, Dense SVD takes leading left singular vectors of the
matrix whose columns are the normalized, aligned task means, following
the use of singular directions in BOLT~\citep{park2026bolt}. If the task means
provide too few directions, we subtract each task mean from its
aligned fits and take leading left singular vectors of the residuals
after projection orthogonal to the means. The statistical basis uses a
single generalized eigenspace per site, pooling the validation samples.
On toroidal MNIST, its signal combines the normalized mean moment at
weight $3/4$ with the normalized paired moment of centered class
covariances at weight $1/4$. Both controls tune their refits on the
common development episodes.

\phantomsection\label{app:shuffle-control}
The shuffled control permutes rows of $\mathbf B_g$ within each
block, preserving group sizes and signs. We recompute the principal
angles to the same $\mathbf U$ and apply the selected rotation fractions.
We average four permutations with matched refit settings and
initialization seeds.

\subsubsection{Frozen-anchor comparisons}

Both comparisons use the Dense source fit with lowest validation loss
as a fixed anchor. On Aug-miniImageNet, diagonal BOLT~\citep{park2026bolt}
initializes 64 source refits from that anchor. For each of 16 patch
matrices, we take one singular direction from each of the first eight
updates, orthogonalize the left and right factors separately and
initialize the diagonal from the projected updates of all 64 fits.
BOLT fits a nonnegative diagonal on each target episode. Both methods
fit fresh hidden biases and output weights. Development selects a
diagonal initialization scale of 10 and 150 Adam steps at rate 0.09
for BOLT, and 300 steps at 0.24 for the fixed anchor.

\begin{table}[htbp]
\centering\small
\caption{Diagonal BOLT on the Aug-miniImageNet episodes of
Table~\ref{tab:transfer-classification}. Both methods transfer the same
fixed Dense anchor. The budget counts adaptable first-layer
coefficients. Preparation time includes the required source fits.}
\label{tab:bolt-diagonal}
\begin{tabular*}{\linewidth}{@{\extracolsep{\fill}}lrrrrr@{}}
\toprule
Method & Coefficients & Clean (\%) & Transformed (\%) & Basis (s) & Preparation (s) \\
\midrule
Fixed Dense anchor & 0 & 32.37 & 29.05 & \textemdash & 21.67 \\
BOLT (diagonal) & 128 & 32.12 & 29.06 & 0.027 & 43.00 \\
\bottomrule
\end{tabular*}
\end{table}

On Aug-Omniglot, the signed-group control parameterizes the update
with the unrotated basis $\mathbf B(\mathbf 0)$ from aSRR. For aTLAS, we fit a Dense
model from the anchor on each of the 48 source tasks and select 32
updates by pivoted QR of their normalized vectors. Each neuron fits
its own mixture at the original update scales. EigenLoRAx fits one
LoRA adapter per source task and applies centered PCA to the input and
hidden factors separately. The selected rank and the number of
principal directions per factor are both 32, giving two $32\times32$
coefficient matrices.

We select the EigenLoRAx dimensions and refits by averaging accuracy
across both views on 192 development episodes. On each test episode,
all three methods start from a zero update and fit its coefficients,
hidden biases and the classifier on five clean training examples per
class. The refits use 75 Adam steps at rates 0.03 for the signed groups, 0.01 for
aTLAS and 0.003 for EigenLoRAx. The paired 95\% intervals for aTLAS
minus signed groups are $[-0.28,0.19]$ points on clean images and $[1.29,1.81]$
on transformed images.

\begin{table}[htbp]
\centering\small
\caption{Test accuracy (\%) with a common frozen Dense anchor on
the 768 Aug-Omniglot episodes of
Table~\ref{tab:transfer-classification}. Each method fits
$2{,}048$ first-layer parameters.}
\label{tab:aug-omniglot-anchored}
\begin{tabular*}{\linewidth}{@{\extracolsep{\fill}}lrr@{}}
\toprule
Method & Clean & Transformed \\
\midrule
\hyperref[alg:signed-merge]{Signed groups with anchor} & \textbf{77.56} & \underline{49.98} \\
\href{https://arxiv.org/abs/2407.02880}{aTLAS}~\citep{zhang2024atlas} & \underline{77.51} & \textbf{51.53} \\
\href{https://arxiv.org/abs/2502.04700}{EigenLoRAx}~\citep{kaushik2025eigenlorax} & 70.91 & 48.21 \\
\bottomrule
\end{tabular*}
\end{table}


\subsection{Compatibility among alternative indexators}
\label{app:relational-compatibility}

Consider a bank of indexators in a common neuron order and let
$V_\tau$ retain all candidates within $\varepsilon$ of the best
validation loss on task $\tau$ after refitting. For an indexator
$\mathbf I$, let $\mathcal R(\mathbf I)$ be its fixed zero
connections and signed equalities. Write $\operatorname{Max}$
for the inclusion-maximal sets of a family. The maximal
compatible sets in Equation~\eqref{eq:method-joint-relations}
can be found by
\[
    \begin{aligned}
        \mathcal A_1&=\operatorname{Max}
            \{\mathcal R(\mathbf I)\mid\mathbf I\in V_1\},\\
        \mathcal A_\tau&=\operatorname{Max}
            \{A\cap\mathcal R(\mathbf I)\mid
               A\in\mathcal A_{\tau-1},\,\mathbf I\in V_\tau\}.
    \end{aligned}
\]
If one set is contained in another, intersecting both with the
same relations preserves inclusion. This justifies pruning at
every step.

In Pattern, a task asks whether an 11-bit input contains a given
4-bit substring. The source bank contains 14 four-code indexators
from Dense fits and source responses. For each of 12 tasks, we fit
on 768 training examples and score the binary Brier loss on 384 validation
examples. With $\varepsilon=0.01$, the recurrence returns one
maximal set. It matches the teacher at all 88 positions up to
neuron permutation, code relabeling and groupwise sign flips.
The teacher has 32 permitted connections and 56 weights fixed at zero.
The result is recovery within the proposal bank.
